\begin{pastexam}[label={px:nchu-115-is-28}]{中興資工所乙組 115 學年度 第 28 題}
（單選；每題 2 分，答錯不倒扣）

Which traversal method is most suitable for detecting whether an undirected graph is connected?

A.~Topological sort \qquad B.~Depth-first search (DFS)\\
C.~Dijkstra's algorithm \qquad D.~Kruskal's algorithm
\end{pastexam}

\begin{pastsoln}
答案是 B。

從任一個頂點做一次 DFS，拜訪到的正是它所在的連通元件（資結書第~21.6~節）。 拜訪到的頂點數
等於 $|V|$ 就是連通，否則不連通。 時間 $O(|V| + |E|)$（鄰接串列）。 BFS 也一樣可以，選項裡只有
DFS。
\begin{itemize}
\item A 錯：拓樸排序只對有向無環圖有意義，無向圖沒有方向，談不上拓樸順序。
\item B 對。
\item C：Dijkstra 從起點出發，最後 $d = \infty$ 的頂點就是走不到的，也能判斷連通；但它要優先佇列，
時間 $O((|V| + |E|)\log |V|)$，還用到不需要的權重，不是最適合的。
\item D：Kruskal 做完時，如果互斥集合只剩一個（加入了 $|V| - 1$ 條邊）就是連通，也判斷得出來；
但它要先把邊排序，$O(|E| \log |E|)$，同樣多做了事。 而且 C、D 都不是「走訪」的方法。
\end{itemize}
\end{pastsoln}
